\documentclass[aspectratio=169]{beamer}

\usepackage{cmap}
\usepackage[T2A]{fontenc}
\usepackage[utf8]{inputenc}
\usepackage[bulgarian]{babel}
\usepackage{tempora}
\usepackage{amsmath, amssymb}
\usepackage{graphicx}
\usepackage{booktabs}
\usepackage{tikz}
\usepackage{hyperref}
\usetikzlibrary{positioning, arrows.meta}
\usetikzlibrary{decorations.pathreplacing}

\usetheme{Copenhagen}
\usecolortheme{default}
\setbeamertemplate{navigation symbols}{}

\title{Моделиране на система за обработка на сензорни данни с крайни срокове}
\subtitle{Курсова работа по „Компютърно моделиране и симулации“}
\author{Алекс Цветанов\\Фак. № 121222225}
\institute{Технически университет -- София\\Факултет „Компютърни системи и технологии“}
\date{}

\begin{document}

\begin{frame}
    \titlepage
\end{frame}

\begin{frame}{Съдържание}
    \tableofcontents
\end{frame}

\section{Постановка на задачата}

\begin{frame}{Описание на системата}
    \begin{itemize}
        \item Разпределена банка от данни с три отдалечени центъра: A, B и C.
        \item Заявките постъпват към централен компютър през интервали \((12 \pm 5)\) минути.
        \item Централният компютър извършва предварителна обработка за \((2 \pm 1)\) минути.
        \item След това се изпращат паралелни заявки към трите центъра.
        \item Обработката приключва след получаване на отговори от всички центрове.
    \end{itemize}
\end{frame}

\begin{frame}{Времеви параметри}
    \begin{table}[ht]
        \centering
        \begin{tabular}{lc}
            \toprule
            Параметър & Стойност \\
            \midrule
            Междупристигащо време & \((12 \pm 5)\) мин \\
            Предварителна обработка & \((2 \pm 1)\) мин \\
            Предаване на заявка & 1 мин \\
            Търсене в център A & \((5 \pm 2)\) мин \\
            Търсене в център B & \((10 \pm 5)\) мин \\
            Търсене в център C & \((8 \pm 4)\) мин \\
            Предаване на отговор & 2 мин \\
            Брой заявки & 150 \\
            \bottomrule
        \end{tabular}
    \end{table}
\end{frame}

\section{Теоретичен модел}

\begin{frame}{Логика на обработката}
    \begin{enumerate}
        \item Генериране на входяща заявка.
        \item Предварителна обработка в централния компютър.
        \item Паралелно изпращане към A, B и C.
        \item Паралелно търсене в трите центъра.
        \item Получаване на трите отговора.
        \item Завършване на заявката след синхронизация.
    \end{enumerate}
\end{frame}

\begin{frame}{Оценка на очакваното време}
    Общото време за една заявка се определя от критичния път:

    \[
        T \approx T_{prep} + T_{send} + \max(T_A, T_B, T_C) + T_{resp}
    \]

    При средни стойности:

    \[
        T \approx 2 + 1 + \max(5, 10, 8) + 2 = 15\ \text{минути}
    \]

    \vspace{0.4cm}
    \begin{itemize}
        \item Очаквано най-силно влияние има център B.
        \item Общото време се доминира от най-бавния паралелен клон.
    \end{itemize}
\end{frame}

\begin{frame}{Архитектура на системата}
    \centering
    \begin{tikzpicture}[node distance=2.6cm, >=Stealth, every node/.style={font=\small}]
        \node[draw, rectangle, align=center] (central) {Централен\\компютър};
        \node[draw, rectangle, align=center, below left=of central] (a) {Център A\\\(5\pm2\) мин};
        \node[draw, rectangle, align=center, below=of central] (b) {Център B\\\(10\pm5\) мин};
        \node[draw, rectangle, align=center, below right=of central] (c) {Център C\\\(8\pm4\) мин};

        \draw[->] (central) to[bend left=15] node[left] {1 мин} (a);
        \draw[->] (central) to[bend left=15] node[right] {1 мин} (b);
        \draw[->] (central) to[bend right=15] node[right] {1 мин} (c);

        \draw[->] (a) to[bend left=15] node[left] {2 мин} (central);
        \draw[->] (b) to[bend left=15] node[left] {2 мин} (central);
        \draw[->] (c) to[bend right=15] node[right] {2 мин} (central);
    \end{tikzpicture}
\end{frame}

\section{Имплементация}

\begin{frame}{Реализация в Python и SimPy}
    \begin{itemize}
        \item Моделът е реализиран на Python с библиотеката \texttt{simpy}.
        \item Входният поток и времената за обработка се генерират с равномерни разпределения.
        \item За паралелното изчакване на трите центъра се използва механизъм от типа \texttt{AllOf}.
        \item За всяка заявка се измерва пълното време от постъпване до приключване.
    \end{itemize}
\end{frame}

\begin{frame}{Ключови параметри на симулацията}
    \begin{block}{Конфигурация}
        \begin{itemize}
            \item \texttt{REQUEST\_COUNT = 150}
            \item \texttt{ARRIVAL\_MEAN = 12}, \texttt{ARRIVAL\_VAR = 5}
            \item \texttt{PREPROCESS\_MIN = 1}, \texttt{PREPROCESS\_MAX = 3}
            \item \texttt{QUERY\_TRANSMIT\_TIME = 1}
            \item \texttt{RESPONSE\_TRANSMIT\_TIME = 2}
            \item \texttt{SEARCH\_A = [3, 7]}, \texttt{SEARCH\_B = [5, 15]}, \texttt{SEARCH\_C = [4, 12]}
        \end{itemize}
    \end{block}
\end{frame}

\begin{frame}{Алтернативна реализация в GPSS}
    \begin{itemize}
        \item Изградена е и независима GPSS реализация за валидация на модела.
        \item Паралелизмът е представен чрез разделяне на транзакцията в отделни клонове.
        \item GPSS е подходящ за проверка на логиката на потока и натоварването.
        \item Python дава по-удобен достъп до агрегирани статистики за времената.
    \end{itemize}
\end{frame}

\begin{frame}[fragile]{GPSS код на модела}
\tiny
\begin{verbatim}
GENERATE 12,5
SEIZE 1
ADVANCE 2,1
RELEASE 1
SPLIT 1,QUERY_A
SPLIT 1,QUERY_B
TRANSFER ,QUERY_C
QUERY_A SEIZE 2
ADVANCE 1
ADVANCE 5,2
ADVANCE 2
RELEASE 2
QUERY_B SEIZE 3
ADVANCE 1
ADVANCE 10,5
ADVANCE 2
RELEASE 3
QUERY_C SEIZE 4
ADVANCE 1
ADVANCE 8,4
ADVANCE 2
RELEASE 4
TERMINATE 0
START 150
\end{verbatim}
\end{frame}

\begin{frame}{Резултати от GPSS симулацията}
    \begin{table}[ht]
        \centering
        \begin{tabular}{lcc}
            \toprule
            Параметър & Стойност \\
            \midrule
            Start time & 0.000 \\
            End time & 1745.583 \\
            Blocks & 37 \\
            Facilities & 6 \\
            Генерирани транзакции & 150 \\
            Завършени транзакции & 150 \\
            \bottomrule
        \end{tabular}
    \end{table}

    \vspace{0.3cm}

    \begin{itemize}
        \item Отчетът показва, че всички 150 заявки са преминали успешно през модела.
        \item GPSS потвърждава коректната синхронизация между трите паралелни клона.
        \item Резултатите са в съответствие с Python симулацията, но без детайлни времеви статистики.
    \end{itemize}
\end{frame}

\section{Резултати}

\begin{frame}{Резултати от симулацията}
    \begin{table}[ht]
        \centering
        \begin{tabular}{lcc}
            \toprule
            Параметър & Python & GPSS \\
            \midrule
            Генерирани заявки & 150 & 150 \\
            Завършени заявки & 150 & 150 \\
            Средно общо време & 20.25 мин & -- \\
            Минимално време & 13.25 мин & -- \\
            Максимално време & 39.36 мин & -- \\
            \bottomrule
        \end{tabular}
    \end{table}
\end{frame}

\begin{frame}{Анализ на резултатите}
    \begin{itemize}
        \item Всички 150 заявки са обработени успешно и в двата модела.
        \item Средното време в Python симулацията е около 20 минути.
        \item Получената средна стойност е по-висока от теоретичната оценка 15 минути.
        \item Причината е вариацията във времената и ефектът на най-бавния паралелен клон.
        \item Най-съществено влияние оказва център B, който има най-голяма средна продължителност.
    \end{itemize}
\end{frame}

\begin{frame}{Сравнение между Python и GPSS}
    \begin{columns}[T]
        \column{0.48\textwidth}
        \begin{block}{Python / SimPy}
            \begin{itemize}
                \item Детайлни временни измервания
                \item Гъвкава статистическа обработка
                \item Лесно разширяване на модела
            \end{itemize}
        \end{block}

        \column{0.48\textwidth}
        \begin{block}{GPSS}
            \begin{itemize}
                \item Ясен транзакционен модел
                \item Подходящ за структурна валидация
                \item По-ограничени обобщени времеви статистики
            \end{itemize}
        \end{block}
    \end{columns}
\end{frame}

\section{Заключение}

\begin{frame}{Основни изводи}
    \begin{enumerate}
        \item Системата работи коректно при паралелна обработка на 150 заявки.
        \item Средното време за обработка е около 20 минути.
        \item Център B определя критичния път в повечето случаи.
        \item Моделът е успешно валидиран чрез Python и GPSS реализация.
    \end{enumerate}
\end{frame}

\begin{frame}{Препоръки за оптимизация}
    \begin{itemize}
        \item Подобряване производителността на център B.
        \item Намаляване на времената за предаване по каналите.
        \item Добавяне на кеширане за често използвани данни.
        \item Очакван ефект: намаляване на средното време с около 25--30\%.
    \end{itemize}
\end{frame}

\begin{frame}
    \centering
    \Large Благодаря за вниманието!
\end{frame}

\end{document}
